\documentclass{article}
\usepackage{amsmath}
\begin{document}

\section{Introduction}
\label{sec:intro}
Background introduction text on embedded systems.

\section{Methodology}
\label{sec:methodology}
The proposed quantized execution engine compiles transformer operators directly to micro-kernels.
The memory layout uses continuous tensor buffers to avoid fragmentation.
Each attention layer computes multi-head projections with specialized SIMD instructions.
Weight quantization uses symmetric per-channel scaling factors calculated during offline calibration.
The activation tensors are dynamically scaled to maintain numerical precision across deep layers.
Intermediate activations are stored in a circular scratchpad buffer located in fast internal memory.
DMA transfers pipeline tensor tiles asynchronously between external flash and internal SRAM.
The execution scheduler dynamically balances compute across vector units and hardware accelerators.
Execution timing is strictly deterministic with zero dynamic memory allocations in the hot loop.
Operator fusion eliminates memory roundtrips between layer normalization and activation functions.

\section{Dependencies}
\label{sec:dependencies}
\subsection{Hardware Architecture}
\label{subsec:hardware_arch}
The MCU features a dual-issue Cortex-M55 core running at 200MHz with Helium vector extensions.
Internal SRAM is partitioned into 1.5MB for instruction and weight cache and 512KB for activation buffers.
External flash operates via Octal-SPI with an effective throughput of 200MB/s for non-cached weights.

\subsection{Quantization Formulae}
\label{subsec:quant_formulae}
Given input tensor $X$, quantized values $q_X$ are computed as $q_X = \text{clamp}(\lfloor X / S \rceil + Z, -128, 127)$.
The scale factor $S = \frac{\max(X) - \min(X)}{255}$ and zero-point $Z = -128 - \lfloor \min(X) / S \rceil$.

\end{document}
